\chapter{Estatística descritiva}\label{ch:g10:stats}

A estatística resume uma grande coleção de números --- notas, salários,
temperaturas --- por meio de alguns indicadores bem escolhidos. Este
capítulo apresenta as duas famílias de indicadores: os que localizam o
\emph{centro} dos dados (\omterm{def:g10:stats:mean}{média}, \omterm{def:g10:stats:median}{mediana}) e os que medem sua
\emph{dispersão} (\omterm{def:g10:stats:quartiles}{amplitude}, quartis). Um tratamento mais profundo,
incluindo o desvio padrão, vem no \cref{ch:g11:stat}.

\section{Dados e frequências}

\begin{definition}[Série estatística]\label{def:g10:stats:series}
Uma \emph{série estatística}\index{série estatística} é uma lista de
valores $x_1, x_2, \dots, x_N$ observados em uma população de tamanho
$N$. Quando um valor $x$ aparece $n$ vezes, $n$ é a sua
\emph{frequência absoluta} e o quociente
\[
f = \frac{n}{N}
\]
é a sua \emph{frequência relativa}\index{frequência relativa} (muitas
vezes expressa em porcentagem). As frequências relativas de todos os
valores somam $1$.
\end{definition}

\begin{example}\label{ex:g10:stats:frequencies}
As notas de $20$ alunos em um teste valendo $5$ pontos:
\begin{center}
\renewcommand{\arraystretch}{1.2}
\begin{tabular}{c|cccccc}
nota $x$ & $0$ & $1$ & $2$ & $3$ & $4$ & $5$ \\
\hline
quantidade $n$ & $1$ & $2$ & $4$ & $6$ & $5$ & $2$ \\
\omterm{def:g10:stats:series}{frequência} $f$ & $0.05$ & $0.10$ & $0.20$ & $0.30$ & $0.25$ & $0.10$ \\
\end{tabular}
\end{center}
Verificação: $1 + 2 + 4 + 6 + 5 + 2 = 20$ e as \omterm{def:g10:stats:series}{frequências} somam $1$.
\end{example}

\begin{omfigure}
\begin{tikzpicture}
\begin{axis}[
  width=8.4cm, height=5.2cm,
  ybar, bar width=14pt,
  xmin=-0.7, xmax=5.7, ymin=0, ymax=7.4,
  xtick={0,1,2,3,4,5},
  ytick={1,2,3,4,5,6},
  xlabel={nota}, ylabel={quantidade},
  tick label style={font=\small},
  label style={font=\small}]
  \addplot[fill=omDef!30, draw=omDef] coordinates
    {(0,1) (1,2) (2,4) (3,6) (4,5) (5,2)};
\end{axis}
\end{tikzpicture}

{\small O \omterm{def:g10:functions:graph}{gráfico} de barras das notas do teste: uma barra por valor, com
altura igual à quantidade.}
\end{omfigure}

\section{Medidas de centro}

\begin{definition}[Média]\label{def:g10:stats:mean}
A \emph{média}\index{média} da série é
\[
\bar x = \frac{x_1 + x_2 + \dots + x_N}{N}.
\]
Se os valores $v_1, \dots, v_k$ aparecem com quantidades
$n_1, \dots, n_k$, o mesmo número é obtido como \emph{média ponderada}
\[
\bar x = \frac{n_1 v_1 + n_2 v_2 + \dots + n_k v_k}{n_1 + n_2 + \dots + n_k}.
\]
\end{definition}

\begin{example}\label{ex:g10:stats:mean}
Para o teste do \cref{ex:g10:stats:frequencies}:
\[
\bar x = \frac{1 \times 0 + 2 \times 1 + 4 \times 2 + 6 \times 3
+ 5 \times 4 + 2 \times 5}{20}
= \frac{0 + 2 + 8 + 18 + 20 + 10}{20}
= \frac{58}{20} = 2.9 .
\]
\end{example}

\begin{definition}[Mediana]\label{def:g10:stats:median}
A \emph{mediana}\index{mediana} é um valor que divide a série
\emph{ordenada} em duas metades: pelo menos metade dos valores é
$\leq$ a mediana e pelo menos metade é $\geq$ a ela. Na prática, ordene
os $N$ valores;
\begin{itemize}
  \item se $N$ é ímpar, a mediana é o valor central, na posição
        $\frac{N+1}{2}$;
  \item se $N$ é par, tome o \omterm{prop:g10:coordgeom:midpoint}{ponto médio} dos dois valores centrais, nas
        posições $\frac N2$ e $\frac N2 + 1$.
\end{itemize}
\end{definition}

\begin{example}\label{ex:g10:stats:median}
Série ordenada de $7$ preços de casas (em milhares): $120$, $150$, $160$,
$180$, $210$, $240$, $900$. A \omterm{def:g10:stats:median}{mediana} é o $4^{\text{o}}$ valor, $180$. A
\omterm{def:g10:stats:mean}{média} é $\frac{1960}{7} = 280$ --- maior do que $6$ dos $7$ preços! Um
valor extremo ($900$) puxa a \omterm{def:g10:stats:mean}{média} muito mais do que a \omterm{def:g10:stats:median}{mediana}: a \omterm{def:g10:stats:median}{mediana}
é \emph{robusta}, a \omterm{def:g10:stats:mean}{média} não.
\end{example}

\section{Medidas de dispersão}

\begin{definition}[Amplitude, quartis]\label{def:g10:stats:quartiles}
Para uma série ordenada:
\begin{itemize}
  \item a \emph{amplitude}\index{amplitude} é a diferença entre o maior e
        o menor valor;
  \item o \emph{primeiro quartil}\index{quartil} $Q_1$ é o menor valor
        tal que pelo menos um quarto ($25\%$) dos valores seja
        $\leq Q_1$; o \emph{terceiro quartil} $Q_3$ é o menor valor tal
        que pelo menos três quartos ($75\%$) dos valores sejam
        $\leq Q_3$;
  \item a \emph{amplitude interquartil}\index{amplitude interquartil}
        $Q_3 - Q_1$ mede a dispersão da metade central dos dados.
\end{itemize}
\end{definition}

\begin{method}[Encontrar os quartis]\label{met:g10:stats:quartiles}
Para uma série de $N$ valores ordenados:
\begin{enumerate}
  \item calcule $\frac N4$; se não for inteiro, arredonde \emph{para
        cima}, até o próximo inteiro; $Q_1$ é o valor nessa posição;
  \item calcule $\frac{3N}{4}$; arredonde para cima se preciso; $Q_3$ é o
        valor nessa posição.
\end{enumerate}
\end{method}

\begin{example}\label{ex:g10:stats:quartiles}
Tome os $11$ valores ordenados
\[
2,\ 3,\ 5,\ 5,\ 6,\ 7,\ 8,\ 9,\ 9,\ 10,\ 12 .
\]
\omterm{def:g10:stats:median}{Mediana}: posição $\frac{11+1}{2} = 6$, logo a \omterm{def:g10:stats:median}{mediana} é $7$.
$Q_1$: $\frac{11}{4} = 2.75$, arredondado para cima dá $3$; o
$3^{\text{o}}$ valor é $5$. $Q_3$: $\frac{33}{4} = 8.25$, arredondado
para cima dá $9$; o $9^{\text{o}}$ valor é $9$. \omterm{def:g10:stats:quartiles}{Amplitude}: $12 - 2 = 10$;
\omterm{def:g10:stats:quartiles}{amplitude interquartil}: $9 - 5 = 4$.
\end{example}

\begin{omfigure}
\begin{tikzpicture}[scale=0.85]
  \draw[-Stealth] (1,0) -- (13.4,0) node[below, font=\small] {valores};
  \foreach \x in {2,4,6,8,10,12}
    { \draw (\x,-0.07) -- (\x,0.07);
      \node[below, font=\footnotesize] at (\x,-0.1) {\x}; }
  \draw[omDef, thick, fill=omDef!12] (5,0.5) rectangle (9,1.3);
  \draw[omDef, very thick] (7,0.5) -- (7,1.3);
  \draw[omDef, thick] (2,0.9) -- (5,0.9);
  \draw[omDef, thick] (9,0.9) -- (12,0.9);
  \draw[omDef, thick] (2,0.7) -- (2,1.1);
  \draw[omDef, thick] (12,0.7) -- (12,1.1);
  \node[above, font=\small] at (5,1.32) {$Q_1$};
  \node[above, font=\small] at (7,1.32) {mediana};
  \node[above, font=\small] at (9,1.32) {$Q_3$};
  \node[above, font=\small] at (2,1.12) {mín};
  \node[above, font=\small] at (12,1.12) {máx};
\end{tikzpicture}

{\small O diagrama de caixa do \cref{ex:g10:stats:quartiles}: a caixa
cobre o \omterm{def:g10:numbers:interval}{intervalo} interquartil $\intcc{Q_1}{Q_3}$ e contém a metade
central dos dados; os bigodes alcançam os valores extremos.}
\end{omfigure}

\begin{example}[Comparar dois grupos]\label{ex:g10:stats:comparison}
Duas turmas fazem a mesma prova (notas de $0$ a $20$):
\begin{center}
\renewcommand{\arraystretch}{1.2}
\begin{tabular}{c|ccc}
 & \omterm{def:g10:stats:median}{mediana} & $Q_1$ & $Q_3$ \\
\hline
turma A & $12$ & $10$ & $14$ \\
turma B & $12$ & $6$ & $17$ \\
\end{tabular}
\end{center}
Mesma \omterm{def:g10:stats:median}{mediana}, dispersões muito diferentes: a turma A é homogênea
(metade de suas notas em $\intcc{10}{14}$), a turma B é heterogênea (sua
metade central cobre $\intcc{6}{17}$). Indicadores de centro sozinhos
nunca contam a história inteira.
\end{example}

\begin{remark}
Nem a \omterm{def:g10:stats:mean}{média} nem a \omterm{def:g10:stats:median}{mediana} é ``melhor'': a \omterm{def:g10:stats:mean}{média} usa todos os valores (e é
necessária para calcular totais), e a \omterm{def:g10:stats:median}{mediana} resiste a valores extremos.
Um resumo sério de uma série apresenta pelo menos um indicador de centro
\emph{e} um de dispersão.
\end{remark}

\section{Exercícios}

\begin{exercise}[$\star$]\label{exo:g10:stats:1}
Eis o número de livros lidos em um ano por $15$ alunos:
\[
0,\ 1,\ 1,\ 2,\ 2,\ 2,\ 3,\ 3,\ 4,\ 4,\ 5,\ 6,\ 7,\ 8,\ 12 .
\]
Calcule a \omterm{def:g10:stats:mean}{média}, a \omterm{def:g10:stats:median}{mediana} e a \omterm{def:g10:stats:quartiles}{amplitude} dessa série.
\end{exercise}

\begin{exercise}[$\star$]\label{exo:g10:stats:2}
Um dado foi lançado $50$ vezes:
\begin{center}
\renewcommand{\arraystretch}{1.2}
\begin{tabular}{c|cccccc}
resultado & $1$ & $2$ & $3$ & $4$ & $5$ & $6$ \\
\hline
quantidade & $7$ & $9$ & $8$ & $10$ & $6$ & $10$ \\
\end{tabular}
\end{center}
Calcule a \omterm{def:g10:stats:series}{frequência} de cada resultado e o resultado médio.
\end{exercise}

\begin{exercise}[$\star$]\label{exo:g10:stats:3}
Encontre a \omterm{def:g10:stats:median}{mediana}, $Q_1$ e $Q_3$ da série ordenada
\[
1,\ 2,\ 4,\ 4,\ 5,\ 6,\ 6,\ 7,\ 8,\ 9,\ 9,\ 11 \qquad (N = 12),
\]
e desenhe o diagrama de caixa correspondente.
\end{exercise}

\begin{exercise}[$\star$]\label{exo:g10:stats:4}
A \omterm{def:g10:stats:mean}{média} de $24$ notas de uma prova é $11.5$. Um $25^{\text{o}}$ aluno faz
a prova e tira $19$. Qual é a nova \omterm{def:g10:stats:mean}{média} da turma?
\end{exercise}

\begin{exercise}[$\star\star$]\label{exo:g10:stats:5}
Um aluno tem \omterm{def:g10:stats:mean}{média} $12$ após $5$ provas, todas com o mesmo peso. Que nota
na $6^{\text{a}}$ prova elevaria a \omterm{def:g10:stats:mean}{média} a $13$? Uma \omterm{def:g10:stats:mean}{média} de $14$ é
alcançável (as notas vão até $20$)?
\end{exercise}

\begin{exercise}[$\star\star$]\label{exo:g10:stats:6}
Invente uma série de $7$ valores não negativos cuja \omterm{def:g10:stats:median}{mediana} seja $10$ e
cuja \omterm{def:g10:stats:mean}{média} seja maior que $20$; depois uma cuja \omterm{def:g10:stats:median}{mediana} seja $10$ e cuja
\omterm{def:g10:stats:mean}{média} seja menor que $6$. Por que a \omterm{def:g10:stats:mean}{média} de uma série assim nunca pode
ser menor que $\frac{40}{7}$? O que esses exemplos mostram?
\end{exercise}

\begin{exercise}[$\star\star$]\label{exo:g10:stats:7}
Em uma empresa, os $9$ funcionários ganham $1800$ cada por mês e o
diretor ganha $12000$.
\begin{enumerate}
  \item Calcule o salário médio e o salário mediano.
  \item Qual indicador descreve melhor um salário ``típico'' aqui? Por
        quê?
\end{enumerate}
\end{exercise}

\begin{exercise}[$\star\star$]\label{exo:g10:stats:8}
Uma turma com $12$ meninos tem altura \omterm{def:g10:stats:mean}{média} de $168$ cm; as $18$ meninas
da mesma turma têm altura \omterm{def:g10:stats:mean}{média} de $163$ cm. Calcule a altura \omterm{def:g10:stats:mean}{média} da
turma inteira. (Cuidado: \emph{não} é o \omterm{prop:g10:coordgeom:midpoint}{ponto médio} entre $168$ e $163$
--- pondere as \omterm{def:g10:stats:mean}{médias} pelos tamanhos dos grupos.)
\end{exercise}

\begin{exercise}[$\star\star\star$]\label{exo:g10:stats:9}
Uma série de $N$ valores $x_1, \dots, x_N$ tem \omterm{def:g10:stats:mean}{média} $\bar x$. Cada valor
é transformado em $y_i = a x_i + b$ para números fixos $a$ e $b$.
\begin{enumerate}
  \item Mostre que a \omterm{def:g10:stats:mean}{média} da nova série é $a \bar x + b$.
  \item As temperaturas de uma semana, em graus Celsius, têm \omterm{def:g10:stats:mean}{média} $20$.
        Qual é a \omterm{def:g10:stats:mean}{média} delas em graus Fahrenheit ($F = 1.8\,C + 32$)?
\end{enumerate}
\end{exercise}

\section{Problema: A pessoa média não existe}

\begin{problem}[{Problema de fim de semana --- média contra mediana:
valores atípicos, curvas de correção, o depósito na estrada e a
cabine projetada para ninguém}]
\label{pb:g10:stats:1}
Em 1950, a Força Aérea dos Estados Unidos mediu $4\,000$ pilotos em
dez dimensões corporais e construiu a cabine para o \emph{homem
médio}. Ela servia --- como um tenente descobriu contando ---
essencialmente a ninguém. A estatística resume multidões com
números isolados, e cada resumo diz a verdade sobre uma coisa e
mente sobre outra. Este problema leva a \omterm{def:g10:stats:mean}{média} e a \omterm{def:g10:stats:median}{mediana}
(\cref{def:g10:stats:mean}, \cref{def:g10:stats:median}) a
julgamento, deixa cada uma mostrar no que é melhor e encerra o
caso da cabine.

\medskip
\noindent\textbf{Parte I --- Dois resumos no banco dos réus.}
\begin{enumerate}
  \item Uma rua vende sete casas, em milhares de euros: $200$,
        $210$, $220$, $230$, $240$, $250$ --- e uma mansão a
        $1\,200$. Calcule o preço médio e o preço mediano. Que
        número descreve ``uma casa desta rua''?
  \item Retire a mansão e recalcule os dois. De quanto cada
        resumo se deslocou? Enuncie a moral sobre valores
        atípicos.
  \item Uma tabela de censo: de $100$ famílias, $20$ não têm
        filhos, $25$ têm um, $30$ têm dois, $15$ têm três e $10$
        têm quatro. Calcule o número médio de filhos. Nenhuma
        família ``tem $1.7$ filhos'' --- então que quantidade real
        a \omterm{def:g10:stats:mean}{média} codifica? (Multiplique-a por $100$.)
  \item Para as quinze notas
        $4, 6, 7, 8, 9, 10, 10, 11, 12, 12, 13, 14, 15, 17, 19$:
        encontre a \omterm{def:g10:stats:median}{mediana} e os quartis $Q_1$, $Q_3$
        (\cref{met:g10:stats:quartiles}).
  \item Para a mesma série, dê a \omterm{def:g10:stats:quartiles}{amplitude} e a \omterm{def:g10:stats:quartiles}{amplitude
        interquartil}. O que cada uma mede?
\end{enumerate}

\medskip
\noindent\textbf{Parte II --- Curvas de correção e turmas
juntadas.}
\begin{enumerate}[resume]
  \item Do \cref{exo:g10:stats:9}: deslocar todos os valores de
        $+b$ desloca \omterm{def:g10:stats:mean}{média}, \omterm{def:g10:stats:median}{mediana} e quartis de $+b$;
        multiplicá-los por $a$ multiplica-os por $a$. O que
        acontece com a \omterm{def:g10:stats:quartiles}{amplitude} e com a \omterm{def:g10:stats:quartiles}{amplitude interquartil}
        sob um deslocamento? E sob uma multiplicação? Justifique.
  \item Uma prova tem \omterm{def:g10:stats:mean}{média} $9$ (de $20$) e o professor quer
        \omterm{def:g10:stats:mean}{média} $12$. Duas correções são propostas: somar $3$
        pontos a todos, ou multiplicar toda nota por $\frac43$.
        Calcule o destino de um aluno com $6$ e de um com $18$
        sob cada correção. Que correção favorece quem --- e qual
        tem um problema de teto?
  \item A turma A ($20$ alunos) tem \omterm{def:g10:stats:mean}{média} $12$; a turma B ($30$
        alunos) tem \omterm{def:g10:stats:mean}{média} $10$. Calcule a \omterm{def:g10:stats:mean}{média} dos $50$ alunos
        juntos e explique por que ela não é $11$ (a lição da
        \omterm{def:g10:stats:mean}{média} ponderada do problema de fim de semana sobre a
        \omterm{def:g10:stats:mean}{média} harmônica, no volume do ensino fundamental, agora
        na sala de aula).
  \item Mostre que as \omterm{def:g10:stats:median}{medianas} recusam esse jogo: calcule as
        \omterm{def:g10:stats:median}{medianas} de $A = \{1, 2, 3\}$ e $B = \{0, 10\}$, depois a
        \omterm{def:g10:stats:median}{mediana} da lista reunida, e compare com a \omterm{def:g10:stats:mean}{média}
        ponderada pelos tamanhos das duas \omterm{def:g10:stats:median}{medianas}.
  \item Patinação artística: cinco juízes dão $5.2$, $5.3$,
        $5.4$, $5.5$, $9.9$. Calcule a \omterm{def:g10:stats:mean}{média}, a \omterm{def:g10:stats:median}{mediana} e a
        \emph{\omterm{def:g10:stats:mean}{média} aparada} (descarte a menor e a maior nota e
        tire a \omterm{def:g10:stats:mean}{média} das restantes). Por que a maioria dos
        esportes com júri apara as notas?
\end{enumerate}

\medskip
\noindent\textbf{Parte III --- Ler distribuições.}
\begin{enumerate}[resume]
  \item Duas fábricas produzem parafusos de comprimento médio
        $50$ mm. Fábrica A: $Q_1 = 49.9$, $Q_3 = 50.1$; fábrica
        B: $Q_1 = 48$, $Q_3 = 52$. Você precisa de parafusos
        intercambiáveis: qual fábrica, e qual estatística
        decidiu?
  \item Em uma clínica, a espera \omterm{def:g10:stats:median}{mediana} é de $10$ minutos, mas a
        espera \omterm{def:g10:stats:mean}{média} é de $25$. Que forma dos dados esse par
        denuncia? Construa um conjunto de cinco valores com
        exatamente essa \omterm{def:g10:stats:median}{mediana} e essa \omterm{def:g10:stats:mean}{média}.
  \item ``Seu bebê está no percentil 75 de altura'' --- o que
        essa frase significa? E todas as crianças de uma cidade
        podem estar ``acima da \omterm{def:g10:stats:mean}{média}''? Acima da \omterm{def:g10:stats:median}{mediana}?
        Explique a diferença.
  \item Turma A: mín $4$, $Q_1\,8$, \omterm{def:g10:stats:median}{mediana} $11$, $Q_3\,13$, máx
        $17$. Turma B: mín $2$, $Q_1\,9$, \omterm{def:g10:stats:median}{mediana} $11$,
        $Q_3\,12$, máx $19$. Compare: nota típica, homogeneidade
        da metade central, selvageria dos extremos.
  \item A lista de verificação do jornalista: uma manchete diz
        ``salário médio na TechCorp: $5\,000$ euros''. Liste as
        três perguntas que este problema lhe ensinou a fazer
        antes de acreditar que a manchete diz alguma coisa sobre
        um funcionário típico.
\end{enumerate}

\medskip
\noindent\textbf{Parte IV --- O caso da cabine.}
\begin{enumerate}[resume]
  \item Construa sua própria prova: monte cinco preços de casas
        com \omterm{def:g10:stats:mean}{média} $300$ e \omterm{def:g10:stats:median}{mediana} $200$ (em milhares de euros).
  \item A \omterm{def:g10:stats:mean}{média} é o ponto de equilíbrio: prove, a partir da
        definição, que os desvios em relação à \omterm{def:g10:stats:mean}{média} sempre se
        cancelam, $\sum (x_i - \bar x) = 0$, e verifique isso no
        seu conjunto de dados da questão 16.
  \item A \omterm{def:g10:stats:median}{mediana} é a que minimiza a rota: há lojas nos
        quilômetros $1$, $3$, $4$, $9$, $13$ de uma estrada, e um
        depósito deve minimizar a \emph{distância total} às cinco
        lojas. Calcule esse total para um depósito na \omterm{def:g10:stats:median}{mediana}
        ($4$), na \omterm{def:g10:stats:mean}{média} ($6$) e em $5$. Quem vence? (Isso é
        geral: a \omterm{def:g10:stats:median}{mediana} minimiza a distância total absoluta ---
        teste qualquer outro ponto.)
  \item A cabine, resolvida: suponha que um piloto tenha
        $30\,\%$ de chance de estar ``perto da \omterm{def:g10:stats:mean}{média}'' em uma
        dimensão corporal, de modo independente ao longo das $10$
        dimensões. Estime a chance de estar perto da \omterm{def:g10:stats:mean}{média} em
        \emph{todas as dez} --- e reconcilie o resultado com o
        que o tenente encontrou entre $4\,000$ pilotos. O que a
        Força Aérea construiu em vez da cabine \omterm{def:g10:stats:mean}{média}?
  \item Final --- o manual de uso, uma linha para cada: quando o
        que importa são totais (um orçamento, uma colheita),
        use\,\dots; quando importa o indivíduo típico (salários,
        tempos de espera), use\,\dots; quando importa a justiça
        da dispersão, informe\,\dots; quando há risco de
        contaminação, use\,\dots. Encerre o caso com o título do
        problema.
\end{enumerate}
\end{problem}
